\section{Mecanismo de ejecución de OS-World: observación, acción y ciclo de interacción}

\subsection{Diseño del espacio de observación}
El curso explicó entradas de observación de múltiples fuentes: capturas de pantalla, el árbol de accesibilidad (accessibility tree), la salida de terminal y flujos de estado personalizados. Cada Agent puede elegir, según sus necesidades, una combinación de entradas dominada por lo visual o por lo textual.

\begin{knowledgebox}{Por qué conviene conservar la captura de pantalla y el árbol de accesibilidad a la vez}
La captura de pantalla ofrece un contexto visual integral, pero con poca estructura; el árbol de accesibilidad cuenta con jerarquía de objetos e información de atributos, lo que facilita la localización precisa y las operaciones interpretables. La complementariedad de ambos mejora de manera notable la estabilidad en la generación de acciones.
\end{knowledgebox}

\begin{figure}[H]
\centering
\includegraphics[width=0.88\textwidth]{frame-03-observation.jpg}
\caption{Video aproximadamente en 00:16:30: el ponente muestra la forma de observación combinada de screenshot y accessibility tree}
\end{figure}

\subsection{Espacio de acción e interfaz de ejecución}
Las acciones se expresan mediante instrucciones ejecutables que abarcan clics, entrada de texto, desplazamiento, atajos de teclado y operaciones combinadas. El Agent no genera directamente una ``sugerencia'', sino código de comportamiento ejecutable en la máquina virtual.

\begin{importantbox}{El valor de estandarizar el formato de las acciones}
Un formato de acción unificado sirve a la vez a tres fines: la alineación de los datos de entrenamiento, la operación estable del ejecutor y la reproducibilidad de los scripts de evaluación. Sin una capa de acciones estándar, tanto la comparación entre modelos como la acumulación de datos a gran escala fracasarían.
\end{importantbox}

\subsection{Ciclo de interacción}
El ciclo del sistema es ``recibir observación $\rightarrow$ planear acción $\rightarrow$ ejecutar en el entorno $\rightarrow$ nueva observación''. La tarea termina al alcanzar el objetivo, al activarse una condición de fallo o al llegar al número máximo de pasos. Este ciclo amplía el objetivo de entrenamiento del Agent, de la predicción de un solo paso a la toma de decisiones secuencial.

\begin{lstlisting}[caption={Ciclo de interacción simplificado de un Agent multimodal en OS-World}]
obs = env.reset(task_config)
for step in range(max_steps):
    action = agent.predict(obs, instruction)
    obs, reward, done, info = env.step(action)
    if done:
        break
\end{lstlisting}

\subsection{Detalles de ingeniería: seguridad y reproducibilidad}
El curso subrayó la importancia del aislamiento en máquinas virtuales y de los archivos de configuración de tareas. El entorno de investigación debe poder reiniciarse, registrarse y reproducirse, de modo que sea posible hacer experimentos comparativos estables y diagnosticar trayectorias de fallo.

\begin{warningbox}{Puntos de fallo comunes}
En un entorno real, lo que más se pasa por alto es la ``deriva de estado'': la caché, el estado de inicio de sesión, las ventanas emergentes y los procesos en segundo plano provocan que la misma tarea se comporte de manera inconsistente entre distintas ejecuciones. Es necesario reducir el efecto de esta deriva mediante instantáneas del entorno, scripts de inicialización y el registro de trazas.
\end{warningbox}

\subsection{Resumen de la sección}
Lo esencial de OS-World no es la captura de pantalla de la interfaz en sí, sino el diseño de ciclo cerrado de ``observación estructurada + acciones estándar + ejecución reproducible''. Esto lleva la investigación de Agents multimodales a una etapa de iteración ingenieril.
